\clearpage
\VerdierPhysicalPageAnchor{0187}%
\begingroup\small\setlength{\parskip}{0pt}\setlength{\emergencystretch}{1.5em}
\setlength{\abovedisplayskip}{4pt plus 2pt minus 2pt}%
\setlength{\belowdisplayskip}{4pt plus 2pt minus 2pt}%

\VerdierSemanticAnchor{ch3d-2-1-1-base-change}
Now let $m\colon Z\longrightarrow X$ be a morphism of complexes. It is
clear that the complex $Z'$ over $Z$ obtained by the base change along
$Z\xrightarrow{m}X$:

\VerdierSemanticAnchor{ch3d-2-1-1-base-change-square}
\begin{center}
\begin{tikzcd}[column sep=4.2em,row sep=2.6em]
Z' \arrow[r] \arrow[d] & X(u,i) \arrow[d,"{p(u,i)}"] \\
Z \arrow[r,"m"'] & X
\end{tikzcd}
\end{center}
is isomorphic to a complex
$p(w,i)\colon Z(w,i)\longrightarrow Z$ for a suitable morphism $w$.

\VerdierSemanticAnchor{ch3d-2-1-2}
\medskip\noindent\textbf{Lemma 2.1.2.}\quad
\textit{a) The morphism
$H^j(p(u,i))\colon H^j(X(u,i))\longrightarrow H^j(X)$ is an
isomorphism for $j\neq i,i+1$.}

\noindent\textit{b) There are exact sequences:}

\VerdierSemanticAnchor{ch3d-2-1-2-exact-sequences}
\[
0\longrightarrow H^i(X(u,i))\longrightarrow H^i(X)
\longrightarrow
\frac{\operatorname{Ker}(d_X^i)+\operatorname{Im}(u)}
 {\operatorname{Im}(d_X^{i-1})+\operatorname{Im}(u)}
\longrightarrow0,
\]
\[
\begin{aligned}
0\longrightarrow{}&
\frac{\operatorname{Ker}(d_X^i)+\operatorname{Im}(u)}
 {\operatorname{Im}(d_X^{i-1})+\operatorname{Im}(u)}
\longrightarrow
\frac{X^i}{\operatorname{Im}(d_X^{i-1})+\operatorname{Im}(u)}\\
&\longrightarrow
\frac{X^i}{\operatorname{Ker}(d_X^i)+\operatorname{Im}(u)}
\longrightarrow0,
\end{aligned}
\]
\[
0\longrightarrow
\frac{X^i}{\operatorname{Ker}(d_X^i)+\operatorname{Im}(u)}
\longrightarrow H^{i+1}(X(u,i))
\longrightarrow H^{i+1}(X)\longrightarrow0.
\]

\noindent\textit{c) The morphism $p(u,i)$ is a quasi-isomorphism if
and only if}
\[
\operatorname{Im}(d_X^{i-1})+\operatorname{Im}(u)=X^i.
\]

\VerdierSemanticAnchor{ch3d-2-1-2-proof}
The diagram chase is left to the reader.

\VerdierSemanticAnchor{ch3d-2-1-3}
\medskip\noindent\textbf{Definition 2.1.3.}\quad
A morphism of complexes $m\colon Z\longrightarrow X$ is called an
\emph{elementary morphism of height $i$} if it is isomorphic in
$\mathcal{F}\ell_1(\operatorname{comp}(\mathcal A))$ to a morphism of the
form $p(u,i)$, where $u$ is an epimorphism.

\VerdierSemanticAnchor{ch3d-2-1-3-properties}
Thus an elementary morphism is, in particular, a quasi-isomorphism
((2.1.2), c)). An elementary morphism is an epimorphism in
$\operatorname{comp}(\mathcal A)$. An epimorphism in
$\operatorname{comp}(\mathcal A)$ is an elementary morphism if and only if
its kernel is an acyclic complex with at most two nonzero components. The
composite of two elementary morphisms of the same height is
\endgroup
